\section{Truncating the Rays}
\label{sec:truncation}

The offsets of a periodic tuple can be large. We use their differences
to find bounded truncations. The cut of a target ray can occur only early
in any admitted comparison. Once its run has rejected, every later
state is a live successor of a rejected state and cannot be the target
of an admitted pair. The host cut can be arbitrarily far
along its ray.

Suppose the onsets are at most $A$ and periods at most $D$. To bound the
target cuts, we use the \emph{horizon} of a ray with onset $a$ and period
$p$, the cut $a+Np$.
Every cut has a representative below this horizon with the same state
and infinite rest. After the onset, the state and the phase determine
all future states and letters. A repetition of a state and phase lets
us move back. Every ray is rejected from some cut on. The representative
of such a cut shows that the first rejected cut lies below
$H=A+ND+1$. An admitted target cut lies at or before the first rejected
cut, so it has index $b<H$.

Let the host cut be $a$ on ray $i$ and the target cut $b$ on ray $j$, with
equal infinite rests. Truncation lengths $T_i$ and $T_j$ avoid this
comparison exactly when the target rest is strictly longer. This is
\eqref{eq:offset} with truncation lengths in place of offsets:
\begin{equation}
   T_i-a < T_j-b .
   \label{eq:strict}
\end{equation}
We impose these difference constraints for every pair of cuts
with equal infinite rests and an admitted pair. This includes $i=j$,
and it includes $a=b$ in that coordinate. For $i=j$ and $a=b$,
\eqref{eq:strict} reads $T_i-a<T_i-a$, which cannot hold. This costs
nothing, because \eqref{eq:offset} already excludes an admitted pair at an
identical cut.

For every needed constraint the lower bound $b-a+1$ on $T_j-T_i$
is at most $H$. The offsets $o_i$ satisfy all these constraints.
Use a directed edge $i\to j$ of weight $\min(H,o_j-o_i)$.
These weights dominate all the required lower bounds and their sum on
a cycle is at most zero. The largest weight of a path ending at each
coordinate, allowing the empty path of weight zero, gives a
nonnegative solution. It is at most $kH$, because repeated vertices
can be removed without decreasing the weight.

Adding the same constant to all these lengths preserves the
differences. We add enough to make any admitted rest of a target ray too
long for a stopped host. If a ray is the host, the matched rest is long
enough that agreement on its letters implies equality of the infinite
rests. Then \eqref{eq:strict} excludes the comparison.

\begin{lemma}[Bounded truncation]\label{lem:truncation}
In a periodic tuple with onsets at most $A$, periods at most $D$,
and stopped word lengths at most $J$, there is a hole of total length at
most
\[
     k(kH+H+A+2D+J),\qquad H=A+ND+1.
\]
\end{lemma}

The proof in \cref{app:truncation} gives the lengths explicitly and
checks every type of comparison. Applying the lemma after the trie
descent gives a bound in the reader parameters.
